\documentclass[10pt,a4paper]{amsart}
\usepackage[margin=19mm]{geometry}
\usepackage{amsmath,amssymb,mathtools}
\usepackage[scr=rsfs,cal=euler]{mathalfa}
\usepackage{xfrac}
\usepackage[unicode,hidelinks,bookmarks=false]{hyperref}
\newcommand{\ie}{i.e.\ }
\newcommand{\cf}{cf.\ }
\newcommand{\resp}{respectively\ }
\newcommand{\Subsection}{\subsection}
\providecommand{\nobreakdash}{\nobreak\hbox{-}}
\setlength{\emergencystretch}{2em}
\multlinegap0pt
\usepackage{xcolor}
\usepackage{algorithm}\usepackage{algpseudocode}
\usepackage{tikz}
\usepackage{amssymb}
\usepackage{mathtools}
\usepackage[shortlabels]{enumitem}
\newtheorem{lemma}{Lemma}
\newtheorem{assumption}{Assumption}
\newtheorem{definition}{Definition}
\title{RanSOM: Second-Order Momentum with Randomized Scaling for Constrained and Unconstrained Optimization}
\author{El Mahdi Chayti}
\date{Source: arXiv:2602.06824v2 (CC BY 4.0)}
\begin{document}
\maketitle

\noindent\emph{Source and licence.} RanSOM: Second-Order Momentum with Randomized Scaling for Constrained and Unconstrained Optimization. Version arXiv:2602.06824v2 (CC BY 4.0), distributed by arXiv under the Creative Commons Attribution 4.0 International licence, http://creativecommons.org/licenses/by/4.0/, as recorded in the arXiv licence metadata for this exact version and shown on the abstract page https://arxiv.org/abs/2602.06824v2. Full licence notice: \texttt{licenses/arxiv-2602.06824-CC-BY-4.0.txt}.\par\smallskip

\noindent\textit{Editorial scope.} Counted unit: the article's lemma lem:app:hess-noise (source0.tex lines 1057--1090). The article's complete original proof of it is delivered with it (source0.tex lines 1104--1387, one \texttt{proof} environment of 284 lines, which the article places away from the statement). No mathematics is written, repaired or re-derived: the statement and the proof are the article's own lines, in the article's own notation, and no retained line is substituted or re-flowed.  Retained uncounted as the source-local context the proof needs: lines 520--533; lines 801--811; lines 813--818; lines 820--836; lines 841--844; lines 964--987; lines 1043--1049.  Nothing was omitted from any retained span, so the package carries no omission record.  No retained line contains a citation, so the wrapper carries no bibliography and no bibliography entry.  No retained line contains a figure environment, an image-inclusion command or drawing code, so no image is delivered and the package declares no asset dependency.  Editorial formatting only, in addition: an AMS \texttt{amsart} preamble that restates the article's own declarations and macros used by the retained lines, namely the environments \texttt{lemma}, \texttt{assumption}, \texttt{definition} on separate counters, together with the standard packages those lines need.  The retained environments print in this document as Lemma 1, Assumption 1, Definition 1 in order of appearance, whereas the article prints the same items with the numbering of its own full document.  The complete native source of the article as distributed in the arXiv e-print for this exact version is bundled unchanged as \texttt{sources/194123/source0.tex}.
\begin{lemma}[Momentum Error Bound]
\label{lemma:main_error}
Under the simplified setting, for any $\beta \in (0, 1/2]$ and $B = 1$,
the averaged momentum error satisfies
\begin{equation}
\label{eq:main_error}
    \frac{1}{T}\sum_{t=0}^{T-1}\mathbb{E}[\|e_t\|_2]
    \le \underbrace{\mathcal{O}\!\left(\frac{\sigma_g}{\beta T\, B_{\mathrm{init}}^{1-1/p}}\right)}_{\text{Initialization}}
    + \underbrace{\mathcal{O}\!\left(\beta^{1-1/p}\,\sigma_g\right)}_{\text{Gradient Noise}}
    + \underbrace{\mathcal{O}\!\left(\eta\,\beta^{-1/q}\,\bar\sigma_h\right)}_{\text{Hessian Bias}},
\end{equation}
where $\bar\sigma_h = \rho(\sigma_h + \kappa L_0)$ is the effective
Hessian-noise constant.
\end{lemma}

\begin{assumption}[Pointwise $(L_0, L_1)$--Smoothness]
\label{ass:app:smooth}
$f$ is twice continuously differentiable and there exist constants
$L_0, L_1 \geq 0$ such that
\begin{equation}
  \|\nabla^2 f(x)\|_{\mathrm{op}}
  \;\leq\; L_0 + L_1 \|\nabla f(x)\|_*
  \qquad \text{for every } x \in \mathcal{E}.
  \label{eq:app:L0L1}
\end{equation}
\end{assumption}

\begin{assumption}[Norm Compatibility]
\label{ass:app:kappa}
There exists $\kappa \geq 1$ such that $\|u\|_* \leq \kappa \|u\|_2$ and
$\|u\|_2 \leq \kappa \|u\|_*$ for every $u \in \mathcal{E}^* \cup
\mathcal{E}$.
\end{assumption}

\begin{assumption}[Affine Heavy--Tailed Noise]
\label{ass:app:noise}
The stochastic gradient $g(x) = \nabla f_\xi(x)$ and the Hessian--vector
product oracle $H(x)w = \nabla^2 f_\xi(x) w$ are unbiased, and there
exist constants $\sigma_g, \alpha_g, \sigma_h, \alpha_h \geq 0$ and
exponents $p, q \in (1, 2]$ such that, for every $x \in \mathcal{E}$ and
every deterministic $w \in \mathcal{E}$ independent of the oracle noise
$\xi$,
\begin{align}
  \mathbb{E}_\xi\bigl[\|g(x) - \nabla f(x)\|_2^p \,\big|\, x\bigr]
    &\leq \sigma_g^p + \alpha_g^p \|\nabla f(x)\|_*^p,
  \label{eq:app:noise-g} \\
  \mathbb{E}_\xi\bigl[\|(H(x) - \nabla^2 f(x))w\|_2^q \,\big|\, x, w\bigr]
    &\leq \bigl(\sigma_h^q + \alpha_h^q \|\nabla f(x)\|_*^q\bigr) \|w\|^q.
  \label{eq:app:noise-h}
\end{align}
\end{assumption}

\begin{equation}
  q L_1 \rho\, \eta_t \;\leq\; \tfrac12,
  \label{eq:app:stepsize-mgf}
\end{equation}

\begin{definition}[Normalized moments]
\label{def:app:moments}
For a given $q \in (1, 2]$, define
\begin{equation}
  M_w \;\triangleq\; \mathbb{E}\!\left[(|w_t|/\eta_t)^q\right],
  \qquad
  M_{ws} \;\triangleq\; \mathbb{E}\!\left[(|w_t|/\eta_t + s_t/\eta_t)^q\right],
  \qquad
  C_s \;\triangleq\; \mathbb{E}[s_t^2]/\eta_t^2.
  \label{eq:app:moments}
\end{equation}
For the unconstrained setting, let $u \triangleq L_1 \rho \eta_t$
satisfying~\eqref{eq:app:stepsize-mgf}, and define
\begin{equation}
  \widetilde C_{B,q} \;\triangleq\;
    \frac{2^{q-1}\bigl(1 + \Gamma(q+1)\bigr)}{(1 - qu)^{q+1}},
  \qquad
  \widetilde C_{A,q} \;\triangleq\;
    2^{q-1}\!\left[M_{ws} +
      u^q\cdot\frac{2^{q-1}\bigl(\Gamma(q+1) + \Gamma(2q+1)\bigr)}
        {(1 - qu)^{2q+1}}\right].
  \label{eq:app:CAB}
\end{equation}
\end{definition}

\begin{equation}
  \Psi_{t+1}
  \;\triangleq\;
  w_t\, H_\xi(x_{t+1})\, d_t
    - \bigl(\nabla f(x_{t+1}) - \nabla f(x_t)\bigr),
  \label{eq:app:Psi-def}
\end{equation}

\begin{lemma}[Hessian--Noise Bound, unconstrained]
\label{lem:app:hess-noise}
Consider RanSOM--E with $s_t \sim \mathrm{Exp}(1/\eta_t)$ and $w_t =
\eta_t$ deterministic. Under
Assumptions~\ref{ass:app:smooth}--\ref{ass:app:noise} and the stepsize
condition~\eqref{eq:app:stepsize-mgf},
\begin{equation}
  \mathbb{E}\bigl[\|\Psi_{t+1}\|_2^q \,\big|\, \mathcal{F}_t\bigr]
  \;\leq\;
  C_\delta^q\, \eta_t^q\,
    \bigl(\bar\sigma_h^q + \bar\alpha_h^q\, \|\nabla f(x_t)\|_*^q\bigr),
  \label{eq:app:hess-noise}
\end{equation}
where
\begin{align}
  \bar\sigma_h^q
  &\triangleq \kappa^q \rho^q\,\bigl[\sigma_h^q
    + L_0^q\, \widetilde C_{A,q}\bigr],
  \label{eq:app:sigma-bar} \\
  \bar\alpha_h^q
  &\triangleq \kappa^q \rho^q\,\bigl[\tfrac{2^{q-1}\,\alpha_h^q}{1 - qu}
    + L_1^q\, \widetilde C_{B,q}\bigr],
  \label{eq:app:alpha-bar} \\
  C_\delta^q
  &\triangleq 2^{2(q-1)}\cdot \max\bigl(M_w,\; M_{ws}\cdot 2^{q-1}\bigr).
  \label{eq:app:C-delta}
\end{align}

Moreover, at $L_1 = 0$ and $\alpha_h = 0$ the formula reduces to
$\bar\sigma_h^q \asymp \kappa^q\rho^q(\sigma_h^q + L_0^q)$, matching
$\bar\sigma_h = \rho(\sigma_h + \kappa L_0)$ of
Lemma~\ref{lemma:main_error} up to a numerical constant absorbed into
$C_\delta$.
\end{lemma}

\begin{proof}[Proof of Lemma~\ref{lem:app:hess-noise}]
We decompose $\Psi_{t+1}$ into a stochastic Hessian--noise term and a
deterministic pathwise approximation term, then bound each separately.

\smallskip\noindent\textbf{Decomposition.}
Writing $x_\tau = x_t + \tau s_t d_t$ for $\tau \in [0, 1]$ (so $x_0 =
x_t$ and $x_1 = x_{t+1}$), the fundamental theorem of calculus gives
\[
  \nabla f(x_{t+1}) - \nabla f(x_t)
  = \int_0^{s_t} \nabla^2 f(x_t + z d_t)\, d_t\, dz
  = s_t \int_0^1 \nabla^2 f(x_\tau)\, d_t\, d\tau.
\]
Substituting into~\eqref{eq:app:Psi-def},
\begin{align}
  \Psi_{t+1}
  &= \underbrace{w_t\bigl[H_\xi(x_{t+1}) - \nabla^2 f(x_{t+1})\bigr] d_t}_{T_1}
  + \underbrace{\Bigl[w_t\, \nabla^2 f(x_{t+1})\, d_t
      - s_t\!\int_0^1 \nabla^2 f(x_\tau)\, d_t\, d\tau\Bigr]}_{T_2}.
  \label{eq:app:Psi-decomp}
\end{align}
By $(a+b)^q \leq 2^{q-1}(a^q + b^q)$ (valid for $a, b \geq 0$ and $q \geq 1$),
\begin{equation}
  \mathbb{E}\bigl[\|\Psi_{t+1}\|_2^q \,\big|\, \mathcal{F}_t\bigr]
  \;\leq\; 2^{q-1}\Bigl(
    \mathbb{E}[\|T_1\|_2^q \mid \mathcal{F}_t]
    + \mathbb{E}[\|T_2\|_2^q \mid \mathcal{F}_t]\Bigr).
  \label{eq:app:Psi-split}
\end{equation}

\smallskip\noindent\textbf{Step 1 (Gradient transfer via Gr\"onwall).}
We first establish a pathwise bound that transfers the gradient norm at
any point $x_\tau$ on the interpolation segment back to the
$\mathcal{F}_t$--measurable anchor $x_t$, using pointwise
$(L_0, L_1)$--smoothness (Assumption~\ref{ass:app:smooth}).

Let $\varphi(\tau) \triangleq \|\nabla f(x_\tau)\|_*$. By the chain rule
and the definition of the induced operator norm $\|\cdot\|_{\mathrm{op}}$,
\[
  \varphi'(\tau)
  = \bigl\langle \tfrac{d}{d\tau}\nabla f(x_\tau),\,
      \mathrm{sign}(\nabla f(x_\tau))\bigr\rangle
  \leq \|\nabla^2 f(x_\tau)\,(s_t d_t)\|_*
  \leq \|\nabla^2 f(x_\tau)\|_{\mathrm{op}}\, s_t \|d_t\|.
\]
(The first inequality is a standard consequence of the fact that
$\varphi$ is $\|\cdot\|_*$--Lipschitz along the path; if one prefers,
$|\varphi(\tau_2) - \varphi(\tau_1)| \leq \|\nabla f(x_{\tau_2}) -
\nabla f(x_{\tau_1})\|_*$ which is $\leq \int$ of the Hessian.) Using
$\|d_t\| \leq \rho$ and Assumption~\ref{ass:app:smooth} pointwise at
$x_\tau$,
\[
  \varphi'(\tau) \leq \rho s_t\, \bigl(L_0 + L_1 \varphi(\tau)\bigr).
\]
This is a linear differential inequality. By Gr\"onwall's lemma (e.g.
multiply both sides by $e^{-L_1\rho s_t \tau}$ and integrate), for all
$\tau \in [0, 1]$,
\begin{equation}
  \varphi(\tau)
  \leq e^{L_1 \rho s_t \tau}\bigl(\varphi(0) + L_0 \rho s_t \tau\bigr)
  \leq e^{L_1 \rho s_t}\bigl(\|\nabla f(x_t)\|_* + L_0 \rho s_t\bigr).
  \label{eq:app:gronwall}
\end{equation}
In particular, at $\tau = 1$,
\begin{equation}
  \|\nabla f(x_{t+1})\|_*
  \leq e^{L_1 \rho s_t}\bigl(\|\nabla f(x_t)\|_* + L_0 \rho s_t\bigr).
  \label{eq:app:gronwall-endpoint}
\end{equation}

\smallskip\noindent\textbf{Step 2 (Bounding $\mathbb{E}[\|T_1\|_2^q \mid \mathcal{F}_t]$).}
The Hessian noise is $T_1 = w_t [H_\xi(x_{t+1}) - \nabla^2 f(x_{t+1})]
d_t$. Condition on everything except the oracle noise:
Assumption~\ref{ass:app:noise} bounds the $\|\cdot\|_2^q$ moment
directly (no norm compatibility needed), with $\|d_t\| \leq \rho$:
\begin{equation}
  \mathbb{E}_\xi\bigl[\|T_1\|_2^q \,\big|\, x_{t+1}, s_t, w_t\bigr]
  \;\leq\; |w_t|^q\, \rho^q\,
    \bigl(\sigma_h^q + \alpha_h^q \|\nabla f(x_{t+1})\|_*^q\bigr).
  \label{eq:app:T1-cond}
\end{equation}
Raise~\eqref{eq:app:gronwall-endpoint} to the $q$-th power and apply
$(a+b)^q \leq 2^{q-1}(a^q + b^q)$:
\begin{equation}
  \|\nabla f(x_{t+1})\|_*^q
  \leq 2^{q-1}\, e^{q L_1 \rho s_t}\,
    \bigl(\|\nabla f(x_t)\|_*^q + L_0^q \rho^q s_t^q\bigr).
  \label{eq:app:grad-transfer-q}
\end{equation}
For RanSOM--E, $w_t = \eta_t$ is deterministic, so $|w_t|^q = \eta_t^q$.
Taking $\mathbb{E}[\,\cdot \mid \mathcal{F}_t]$ of~\eqref{eq:app:T1-cond}
and substituting~\eqref{eq:app:grad-transfer-q}:
\begin{align}
  \mathbb{E}[\|T_1\|_2^q \mid \mathcal{F}_t]
  &\leq \eta_t^q\rho^q\,\sigma_h^q
    + 2^{q-1}\eta_t^q \rho^q\alpha_h^q\,
      \mathbb{E}\!\left[e^{qL_1\rho s_t}\right] \|\nabla f(x_t)\|_*^q
  \notag \\
  &\quad + 2^{q-1}\eta_t^q \rho^q\alpha_h^q L_0^q \rho^q\,
      \mathbb{E}\!\left[s_t^q e^{qL_1\rho s_t}\right],
  \label{eq:app:T1-expanded}
\end{align}
where we used that $\|\nabla f(x_t)\|_*$ is $\mathcal{F}_t$--measurable.

The two exponential--type moments admit closed forms. For $s_t \sim
\mathrm{Exp}(1/\eta_t)$ and any $\lambda < 1/\eta_t$,
\begin{equation}
  \mathbb{E}\!\left[s_t^k e^{\lambda s_t}\right]
  = \frac{1}{\eta_t}\int_0^\infty z^k e^{-z(1/\eta_t - \lambda)}\,dz
  = \frac{\Gamma(k+1)\,\eta_t^k}{(1 - \lambda \eta_t)^{k+1}}.
  \label{eq:app:exp-moment}
\end{equation}
Setting $\lambda = qL_1\rho$ and $u = L_1\rho\eta_t$ (so $\lambda\eta_t =
qu \leq 1/2$ by~\eqref{eq:app:stepsize-mgf}):
\begin{equation}
  \mathbb{E}[e^{qL_1\rho s_t}] = \frac{1}{1 - qu}, \qquad
  \mathbb{E}[s_t^q e^{qL_1\rho s_t}]
    = \frac{\Gamma(q+1)\, \eta_t^q}{(1 - qu)^{q+1}}.
  \label{eq:app:exp-mgf}
\end{equation}
Substituting into~\eqref{eq:app:T1-expanded},
\begin{equation}
  \mathbb{E}[\|T_1\|_2^q \mid \mathcal{F}_t]
  \leq \eta_t^q \rho^q \!\left[\sigma_h^q
    + \frac{2^{q-1}\alpha_h^q}{1 - qu}\, \|\nabla f(x_t)\|_*^q
    + \alpha_h^q L_0^q \rho^q \eta_t^q\,
      \frac{2^{q-1}\Gamma(q+1)}{(1-qu)^{q+1}}\right].
  \label{eq:app:T1-final}
\end{equation}

\smallskip\noindent\textbf{Step 3 (Pathwise bound for $T_2$).}
By definition,
\[
  T_2 = w_t\, \nabla^2 f(x_{t+1})\, d_t
    - s_t \int_0^1 \nabla^2 f(x_\tau)\, d_t\, d\tau.
\]
Applying the triangle inequality and the induced--operator--norm
inequality $\|\nabla^2 f(x) u\|_* \leq \|\nabla^2 f(x)\|_{\mathrm{op}}
\|u\|$, then converting $\|\cdot\|_*$ to $\|\cdot\|_2$ via
Assumption~\ref{ass:app:kappa}:
\begin{equation}
  \|T_2\|_2 \leq \kappa\, \|T_2\|_*
  \leq \kappa\rho\Bigl(|w_t|\cdot \|\nabla^2 f(x_{t+1})\|_{\mathrm{op}}
    + s_t \int_0^1 \|\nabla^2 f(x_\tau)\|_{\mathrm{op}}\, d\tau\Bigr).
  \label{eq:app:T2-norm}
\end{equation}
Apply Assumption~\ref{ass:app:smooth} pointwise at each $x_\tau$ and
$x_{t+1}$, then use~\eqref{eq:app:gronwall}:
\[
  \|\nabla^2 f(x_\tau)\|_{\mathrm{op}}
  \leq L_0 + L_1 \varphi(\tau)
  \leq L_0 + L_1 e^{L_1\rho s_t}\bigl(\|\nabla f(x_t)\|_* + L_0\rho s_t\bigr),
\]
uniform in $\tau \in [0,1]$; the same bound holds at $\tau = 1$ for
$\|\nabla^2 f(x_{t+1})\|_{\mathrm{op}}$. Substituting
into~\eqref{eq:app:T2-norm},
\begin{equation}
  \|T_2\|_2
  \leq \kappa\rho\, (|w_t| + s_t)\,\Bigl[L_0
    + L_1 e^{L_1\rho s_t}\bigl(\|\nabla f(x_t)\|_* + L_0 \rho s_t\bigr)\Bigr].
  \label{eq:app:T2-pathwise}
\end{equation}

\smallskip\noindent\textbf{Step 4 (Splitting $T_2$ into baseline and gradient--linear parts).}
Distribute the factor $(|w_t| + s_t)$ over the bracket in
\eqref{eq:app:T2-pathwise} and collect by powers of $\|\nabla
f(x_t)\|_*$:
\[
  \|T_2\|_2 \leq \kappa\rho\,\bigl(\widetilde A
    + \widetilde B\, \|\nabla f(x_t)\|_*\bigr),
\]
where
\begin{align}
  \widetilde A &:= L_0\, (|w_t| + s_t)\,
    \bigl(1 + L_1\rho s_t e^{L_1\rho s_t}\bigr),
  \label{eq:app:A-tilde} \\
  \widetilde B &:= L_1\, (|w_t| + s_t)\, e^{L_1\rho s_t}.
  \label{eq:app:B-tilde}
\end{align}
Since $(s_t, w_t)$ are independent of $\mathcal{F}_t$ and $\|\nabla
f(x_t)\|_*$ is $\mathcal{F}_t$--measurable, $(a+b)^q \leq 2^{q-1}(a^q +
b^q)$ gives
\begin{equation}
  \mathbb{E}[\|T_2\|_2^q \mid \mathcal{F}_t]
  \leq \kappa^q \rho^q\, 2^{q-1}
    \bigl(\mathbb{E}[\widetilde A^q]
      + \mathbb{E}[\widetilde B^q]\, \|\nabla f(x_t)\|_*^q\bigr).
  \label{eq:app:T2-cond}
\end{equation}

\smallskip\noindent\textbf{Step 5 (Moments of $\widetilde B$).}
By $(a+b)^q \leq 2^{q-1}(a^q + b^q)$,
\begin{align*}
  \mathbb{E}[\widetilde B^q]
  &= L_1^q\, \mathbb{E}\!\left[(|w_t|+s_t)^q e^{qL_1\rho s_t}\right]
  \\
  &\leq L_1^q\, 2^{q-1}\!\left(|\eta_t|^q\, \mathbb{E}[e^{qL_1\rho s_t}]
      + \mathbb{E}[s_t^q e^{qL_1\rho s_t}]\right)
  \qquad (\text{using } w_t = \eta_t).
\end{align*}
Substituting~\eqref{eq:app:exp-mgf},
\[
  \mathbb{E}[\widetilde B^q]
  \leq L_1^q\, 2^{q-1}\!\left[\frac{\eta_t^q}{1-qu}
    + \frac{\Gamma(q+1)\, \eta_t^q}{(1-qu)^{q+1}}\right]
  \leq L_1^q\, \eta_t^q\cdot
    \frac{2^{q-1}\bigl(1 + \Gamma(q+1)\bigr)}{(1-qu)^{q+1}}
  = L_1^q\, \eta_t^q\, \widetilde C_{B,q},
\]
where we used $(1-qu)^{-1} \leq (1-qu)^{-(q+1)}$ since $0 \leq qu < 1$.

\smallskip\noindent\textbf{Step 6 (Moments of $\widetilde A$).}
Applying $(a+b)^q \leq 2^{q-1}(a^q + b^q)$ inside
\eqref{eq:app:A-tilde}:
\[
  \widetilde A^q
  \leq L_0^q\, 2^{q-1}\bigl[(|w_t|+s_t)^q
    + (L_1\rho)^q\, (|w_t|+s_t)^q\, s_t^q\, e^{qL_1\rho s_t}\bigr].
\]
The first term's expectation is at most $M_{ws}\eta_t^q$ by
Definition~\ref{def:app:moments}. For the second, apply $(a+b)^q \leq
2^{q-1}(a^q+b^q)$ once more to $(|w_t|+s_t)^q = (\eta_t + s_t)^q \leq
2^{q-1}(\eta_t^q + s_t^q)$:
\begin{align*}
  \mathbb{E}\!\left[(\eta_t + s_t)^q s_t^q e^{qL_1\rho s_t}\right]
  &\leq 2^{q-1}\bigl(\eta_t^q\, \mathbb{E}[s_t^q e^{qL_1\rho s_t}]
    + \mathbb{E}[s_t^{2q} e^{qL_1\rho s_t}]\bigr)
  \\
  &\leq 2^{q-1}\!\left(\frac{\Gamma(q+1)\,\eta_t^{2q}}{(1-qu)^{q+1}}
    + \frac{\Gamma(2q+1)\,\eta_t^{2q}}{(1-qu)^{2q+1}}\right)
  \\
  &\leq \eta_t^{2q}\cdot
    \frac{2^{q-1}\bigl(\Gamma(q+1) + \Gamma(2q+1)\bigr)}{(1-qu)^{2q+1}},
\end{align*}
using $(1-qu)^{-(q+1)} \leq (1-qu)^{-(2q+1)}$. Multiplying by $(L_1\rho)^q$
and noting $(L_1\rho\eta_t)^q = u^q$,
\[
  (L_1\rho)^q\, \mathbb{E}\!\left[(\eta_t + s_t)^q s_t^q e^{qL_1\rho s_t}\right]
  \leq \eta_t^q \cdot u^q \cdot
    \frac{2^{q-1}\bigl(\Gamma(q+1) + \Gamma(2q+1)\bigr)}{(1-qu)^{2q+1}}.
\]
Combining, $\mathbb{E}[\widetilde A^q] \leq L_0^q\, \eta_t^q\, \widetilde
C_{A,q}$ with $\widetilde C_{A,q}$ as defined in~\eqref{eq:app:CAB}. The
$u^q$ factor is bounded by $2^{-q}$ under~\eqref{eq:app:stepsize-mgf}, so
the second term in the bracket of $\widetilde C_{A,q}$ is a numerical
constant.

\smallskip\noindent\textbf{Step 7 (Assembly).}
Substituting the moments of $\widetilde A$ and $\widetilde B$
into~\eqref{eq:app:T2-cond},
\begin{equation}
  \mathbb{E}[\|T_2\|_2^q \mid \mathcal{F}_t]
  \leq \kappa^q \rho^q\, 2^{q-1}\, \eta_t^q\,
    \bigl(L_0^q\, \widetilde C_{A,q}
      + L_1^q\, \widetilde C_{B,q}\, \|\nabla f(x_t)\|_*^q\bigr).
  \label{eq:app:T2-final}
\end{equation}
Now combine~\eqref{eq:app:T1-final} and~\eqref{eq:app:T2-final}
via~\eqref{eq:app:Psi-split}. The baseline $\sigma_h$--term of $T_1$
gives $\eta_t^q\rho^q\sigma_h^q$, multiplied by the outer $2^{q-1}$
from~\eqref{eq:app:Psi-split}; the $\widetilde A$--term of $T_2$ gives
$\kappa^q\rho^q 2^{q-1}\cdot 2^{q-1} L_0^q\widetilde C_{A,q}\eta_t^q =
\kappa^q\rho^q 2^{2(q-1)} L_0^q \widetilde C_{A,q}\eta_t^q$. Their sum,
together with the factor $M_w = 1$ for RanSOM--E, is absorbed into
$C_\delta^q\eta_t^q\bar\sigma_h^q$ with $\bar\sigma_h$ as
in~\eqref{eq:app:sigma-bar} and $C_\delta^q$ as
in~\eqref{eq:app:C-delta}.

The gradient--coefficient terms: from $T_1$, $2^{q-1}\alpha_h^q/(1-qu)$;
from $T_2$, $2^{2(q-1)} L_1^q\widetilde C_{B,q}$. Their sum
is~\eqref{eq:app:alpha-bar} times $\kappa^q\rho^q$, absorbed into
$C_\delta^q\bar\alpha_h^q$.

The higher--order term $\alpha_h^q L_0^q\rho^q\eta_t^q \cdot
2^{q-1}\Gamma(q+1)/(1-qu)^{q+1}$ from~\eqref{eq:app:T1-final} carries an
extra $\eta_t^q$; since $\eta_t \leq \bar\eta$ (finite), it may be
absorbed into $\bar\sigma_h^q$ with a bounded multiplicative constant.

This proves~\eqref{eq:app:hess-noise}. At $L_1 = 0$ and $\alpha_h = 0$,
$\widetilde C_{A,q}$ reduces to $2^{q-1} M_{ws}$, so $\bar\sigma_h^q =
\kappa^q\rho^q[\sigma_h^q + 2^{q-1}M_{ws} L_0^q] \asymp
\kappa^q\rho^q(\sigma_h + L_0)^q$, matching $\bar\sigma_h = \rho(\sigma_h
+ \kappa L_0)$ of the main text's Lemma~\ref{lemma:main_error} up to a
numerical constant absorbed into $C_\delta$.
\end{proof}

\end{document}
